%% Journal of Great Lakes Research (Elsevier) -- elsarticle, author-year style
%% Plantilla oficial: elsarticle.cls + elsarticle-harv.bst (CTAN, ene-2026)
%%
%% Los numeros de este manuscrito NO se teclean a mano: se importan de
%% numbers.tex, que genera pipeline/p09_export_manuscript_numbers.py a partir
%% de los artefactos del analisis. Regenerar con:
%%     python pipeline/p09_export_manuscript_numbers.py
\documentclass[preprint,review,12pt,authoryear]{elsarticle}

\usepackage{graphicx}
\usepackage{amsmath,amssymb}
\usepackage{booktabs}
\usepackage{multirow}
\usepackage{url}
%% Permite partir URLs largas en cualquier caracter, para que el enlace del
%% repositorio no se salga del margen.
\usepackage{xurl}
\usepackage{lineno}
\linenumbers
\graphicspath{{figures/}{../../results/figures_r/}}

%% Las tablas de este articulo son anchas; \tabfit las ajusta a la caja de
%% texto sin recurrir a \resizebox, que deformaria el tipo.
\newcommand{\tabfit}{\small\setlength{\tabcolsep}{4pt}\renewcommand{\arraystretch}{1.08}}

% numbers.tex -- generado por pipeline/p09_export_manuscript_numbers.py
% NO EDITAR A MANO. Regenerar con:  python pipeline/p09_export_manuscript_numbers.py
% Uso en el texto:  \numHeadlineRtwo  en lugar de teclear 0.574

\newcommand{\numBestModelLabel}{Ensemble (RF+ET+XGB+LGBM optimizados)}
\newcommand{\numBestModelRTwo}{0.583}
\newcommand{\numCampaignDatesMedian}{13}
\newcommand{\numCampaignN}{812}
\newcommand{\numCampaignRTwo}{0.574}
\newcommand{\numCampaignRmse}{1.86}
\newcommand{\numCampaignSpanMaxDays}{27}
\newcommand{\numCampaignSpanMinDays}{7}
\newcommand{\numCampaigns}{120}
\newcommand{\numClassicalBoundedRTwo}{0.097}
\newcommand{\numClassicalBoundedRmse}{2.71}
\newcommand{\numClassicalBrokenMaxpred}{143.5}
\newcommand{\numClassicalBrokenRTwo}{-2.458}
\newcommand{\numClassicalBrokenRmse}{5.29}
\newcommand{\numConformalCovFiveZero}{50.9}
\newcommand{\numConformalCovEightZero}{79.1}
\newcommand{\numConformalCovNineZero}{89.4}
\newcommand{\numConformalCovNineFive}{94.3}
\newcommand{\numConformalCovClearest}{71.3}
\newcommand{\numConformalCovTurbidest}{79.5}
\newcommand{\numConformalMaxError}{0.94}
\newcommand{\numConformalWidthFiveZero}{2.50}
\newcommand{\numConformalWidthEightZero}{4.61}
\newcommand{\numConformalWidthNineZero}{6.12}
\newcommand{\numConformalWidthNineFive}{7.37}
\newcommand{\numCrossSensorPairs}{272}
\newcommand{\numDateVsYearGapRTwo}{0.079}
\newcommand{\numFieldCampaignsN}{9}
\newcommand{\numHeadlineBias}{-0.02}
\newcommand{\numHeadlineMae}{1.44}
\newcommand{\numHeadlineN}{812}
\newcommand{\numHeadlineRTwo}{0.574}
\newcommand{\numHeadlineRmse}{1.86}
\newcommand{\numIndependentEvents}{540}
\newcommand{\numLinearRTwo}{0.323}
\newcommand{\numLinearRmse}{2.34}
\newcommand{\numLocalRBlue}{0.26}
\newcommand{\numLocalRGreen}{0.71}
\newcommand{\numLocalRNir}{0.62}
\newcommand{\numLocalRRed}{0.52}
\newcommand{\numLocalRSwirOne}{0.64}
\newcommand{\numLocalRSwirTwo}{0.60}
\newcommand{\numLocalReharmN}{812}
\newcommand{\numLocalReharmRTwo}{0.518}
\newcommand{\numLocalReharmRmse}{1.98}
\newcommand{\numLsOnlyN}{538}
\newcommand{\numLsOnlyRTwo}{0.549}
\newcommand{\numLsOnlyRmse}{1.81}
\newcommand{\numLsOverSTwoBlue}{1.7}
\newcommand{\numLsOverSTwoGreen}{1.6}
\newcommand{\numLsOverSTwoNir}{15.4}
\newcommand{\numLsOverSTwoRed}{2.7}
\newcommand{\numLsOverSTwoSwirOne}{13.3}
\newcommand{\numLsOverSTwoSwirTwo}{5.7}
\newcommand{\numMatchupsLs}{662}
\newcommand{\numMatchupsSTwo}{340}
\newcommand{\numMatchupsSecchi}{812}
\newcommand{\numMatchupsTotal}{1002}
\newcommand{\numMonthFirst}{7}
\newcommand{\numMonthLast}{10}
\newcommand{\numNullCampaignDateRTwo}{-0.012}
\newcommand{\numNullStationRTwo}{0.208}
\newcommand{\numNullZonaRTwo}{0.287}
\newcommand{\numNullgridCampaignDateCampaignDateRTwo}{-0.012}
\newcommand{\numNullgridCampaignDateStationRTwo}{0.208}
\newcommand{\numNullgridCampaignDateZonaRTwo}{0.287}
\newcommand{\numNullgridStationCampaignDateRTwo}{0.305}
\newcommand{\numNullgridStationStationRTwo}{-0.008}
\newcommand{\numNullgridStationZonaRTwo}{0.276}
\newcommand{\numNullgridYearCampaignDateRTwo}{-0.032}
\newcommand{\numNullgridYearStationRTwo}{0.222}
\newcommand{\numNullgridYearZonaRTwo}{0.234}
\newcommand{\numOneEventN}{540}
\newcommand{\numOneEventRTwo}{0.571}
\newcommand{\numOneEventRmse}{1.77}
\newcommand{\numPairedEvents}{270}
\newcommand{\numRandomN}{812}
\newcommand{\numRandomRTwo}{0.596}
\newcommand{\numRandomRmse}{1.81}
\newcommand{\numRatiosN}{812}
\newcommand{\numRatiosRTwo}{0.468}
\newcommand{\numRatiosRmse}{2.08}
\newcommand{\numRetrChlN}{1002}
\newcommand{\numRetrChlRTwo}{0.012}
\newcommand{\numRetrLstThermalN}{662}
\newcommand{\numRetrLstThermalRTwo}{0.020}
\newcommand{\numRetrSecchiN}{812}
\newcommand{\numRetrSecchiRTwo}{0.574}
\newcommand{\numRetrTempInsituN}{1002}
\newcommand{\numRetrTempInsituRTwo}{0.113}
\newcommand{\numRetrTssN}{920}
\newcommand{\numRetrTssRTwo}{0.056}
\newcommand{\numRfRTwo}{0.574}
\newcommand{\numRfRmse}{1.86}
\newcommand{\numRoyInterceptMax}{0.0448}
\newcommand{\numRoyInterceptMin}{0.0116}
\newcommand{\numRoyRatioBlue}{2.5}
\newcommand{\numRoyRatioGreen}{1.4}
\newcommand{\numRoyRatioMax}{69.8}
\newcommand{\numRoyRatioNir}{69.8}
\newcommand{\numRoyRatioRed}{10.6}
\newcommand{\numRoyRatioSwirOne}{23.2}
\newcommand{\numRoyRatioSwirTwo}{8.2}
\newcommand{\numSTwoOnlyN}{274}
\newcommand{\numSTwoOnlyRTwo}{0.563}
\newcommand{\numSTwoOnlyRmse}{2.04}
\newcommand{\numSecchiMax}{16.5}
\newcommand{\numSecchiMean}{9.1}
\newcommand{\numSecchiMin}{1.5}
\newcommand{\numSecchiSd}{2.8}
\newcommand{\numSecchiUniqueValues}{73}
\newcommand{\numSeedNoiseRange}{0.003}
\newcommand{\numSeedNoiseSd}{0.001}
\newcommand{\numSensorAccAll}{99.8}
\newcommand{\numSensorAccChance}{66.3}
\newcommand{\numSensorAccRatios}{99.5}
\newcommand{\numSensorAccVisible}{97.9}
\newcommand{\numShapShareNirSwir}{14.4}
\newcommand{\numShapShareOtras}{13.0}
\newcommand{\numShapShareVerdeAzul}{72.6}
\newcommand{\numStationN}{812}
\newcommand{\numStationRTwo}{0.596}
\newcommand{\numStationRmse}{1.81}
\newcommand{\numStations}{143}
\newcommand{\numTemporalCut}{2019}
\newcommand{\numTemporalN}{226}
\newcommand{\numTemporalOverlapPct}{82.6}
\newcommand{\numTemporalRTwo}{0.507}
\newcommand{\numTemporalRmse}{2.63}
\newcommand{\numTemporalShared}{57}
\newcommand{\numTemporalTestStations}{69}
\newcommand{\numTrendBahiaDePunoP}{0.241}
\newcommand{\numTrendBahiaDePunoSlope}{0.16}
\newcommand{\numTrendCombinations}{12}
\newcommand{\numTrendFlips}{5}
\newcommand{\numTrendLagoMayorP}{0.138}
\newcommand{\numTrendLagoMayorSlope}{0.25}
\newcommand{\numTrendLagoMenorP}{1.000}
\newcommand{\numTrendLagoMenorSlope}{0.00}
\newcommand{\numTuningGainRTwo}{0.009}
\newcommand{\numVisibleN}{812}
\newcommand{\numVisibleRTwo}{0.541}
\newcommand{\numVisibleRmse}{1.93}
\newcommand{\numYearFirst}{2013}
\newcommand{\numYearLast}{2024}
\newcommand{\numYearN}{812}
\newcommand{\numYearRTwo}{0.495}
\newcommand{\numYearRmse}{2.02}
\newcommand{\numYearsMissing}{2020, 2021, 2023}
\newcommand{\numZoneBahiaDePunoBias}{1.01}
\newcommand{\numZoneBahiaDePunoN}{88}
\newcommand{\numZoneBahiaDePunoRTwo}{0.221}
\newcommand{\numZoneBahiaDePunoRmse}{1.90}
\newcommand{\numZoneLagoMayorBias}{-0.19}
\newcommand{\numZoneLagoMayorN}{591}
\newcommand{\numZoneLagoMayorRTwo}{0.448}
\newcommand{\numZoneLagoMayorRmse}{1.83}
\newcommand{\numZoneLagoMenorBias}{0.08}
\newcommand{\numZoneLagoMenorN}{133}
\newcommand{\numZoneLagoMenorRTwo}{0.211}
\newcommand{\numZoneLagoMenorRmse}{1.94}
\newcommand{\numZoneoutMayorN}{591}
\newcommand{\numZoneoutMayorRTwo}{-0.047}
\newcommand{\numZoneoutMayorRmse}{2.52}
\newcommand{\numZoneoutMenorN}{133}
\newcommand{\numZoneoutMenorRTwo}{0.247}
\newcommand{\numZoneoutMenorRmse}{1.90}
\newcommand{\numZoneoutPunoN}{88}
\newcommand{\numZoneoutPunoRTwo}{0.036}
\newcommand{\numZoneoutPunoRmse}{2.11}


\journal{Journal of Great Lakes Research}

\begin{document}

\begin{frontmatter}

\title{Water transparency retrieval in a high-altitude large lake from
Sentinel-2 and Landsat 8/9: cross-sensor harmonization limits, validation
design, and calibrated uncertainty in Lake Titicaca}

\author[unap]{Dina Maribel Yana-Yucra}
\author[unap]{Richar Andre Vilca-Solorzano\corref{cor1}}
\ead{75521963@est.unap.edu.pe}
\author[unap]{Fred Torres-Cruz}
\cortext[cor1]{Corresponding author.}
\affiliation[unap]{organization={Faculty of Statistical and Computer
Engineering, Universidad Nacional del Altiplano (UNAP)},
  addressline={Av. Floral 1153},
  city={Puno},
  postcode={21001},
  country={Peru}}

\begin{abstract}
Monitoring water transparency in high-altitude oligotrophic lakes is critical for freshwater management, yet it remains constrained by the sparsity of field data. While multi-mission satellite archives (e.g., harmonized Sentinel-2 and Landsat 8/9) and machine learning offer a synoptic solution, current practices often import assumptions from turbid-water literature that fail in clear-water, campaign-based settings. We assemble \numMatchupsTotal{} satellite--field match-ups against twelve years of limnological monitoring in Lake Titicaca to address these failures and define a defensible retrieval protocol. First, we demonstrate that standard land-derived harmonization coefficients (e.g., \citealp{Roy2016}) are invalid over dark water, as their intercepts exceed the native Landsat signal by factors of \numRoyRatioGreen{} to \numRoyRatioNir{}, leaving sensors \numSensorAccAll\% separable; we propose a classifier two-sample test to catch this failure. Second, we show that standard station-blocked cross-validation fails to prevent spatial leakage in campaign-based sampling: under campaign-blocked folds, a station-mean null model still explains $R^2=\numNullStationRTwo$ of the variance. We introduce a null-model diagnostic to identify the correct blocking unit (the campaign). Finally, under this strict design, we train a Random Forest regressor to retrieve Secchi disk depth ($R^2=\numHeadlineRTwo$, RMSE \numHeadlineRmse~m) and attach split-conformal prediction intervals. We reveal that while marginally calibrated (\numConformalCovNineZero\% coverage at 90\% nominal), conformal intervals structurally under-cover in the clearest water (\numConformalCovClearest\%). Our findings provide a rigorous blueprint for evaluating machine-learning water quality products in large, oligotrophic lakes.
\end{abstract}

\begin{highlights}
\item Land-derived Landsat/Sentinel-2 harmonization is invalid over dark
oligotrophic water
\item Sensors stay 99.8\% separable after harmonization; intercepts exceed
signal 70-fold
\item Campaign, not station, is the blocking unit in campaign-based
monitoring programmes
\item Honest Secchi retrieval: $R^2$=0.57, RMSE 1.86 m over 812 real
match-ups
\item Conformal intervals are calibrated marginally but under-cover in the
clearest water
\end{highlights}

\begin{keyword}
Lake Titicaca \sep Secchi disk depth \sep Sentinel-2 \sep Landsat \sep
cross-sensor harmonization \sep conformal prediction \sep large lakes
\end{keyword}

\end{frontmatter}

\section{Introduction}\label{sec:intro}

Large lakes are sentinels of environmental change, but continuous monitoring of remote, high-altitude systems is constrained by sparse field infrastructure \citep{Williamson2009}. Lake Titicaca (3,810~m a.s.l.), the largest tropical lake in the world by volume, is a critical transboundary freshwater reserve whose water quality is documented only by infrequent field campaigns \citep{OEFA2016,Siguayro2022eval}, even as eutrophication advances in its near-shore bays \citep{Heredia2022,Lanza2024}. Satellite remote sensing offers a synoptic solution \citep{Gholizadeh2016,Topp2020}. Specifically, water transparency---quantified by the Secchi disk depth ($Z_{SD}$)---is an integrative indicator of water clarity strongly linked to light attenuation \citep{Lee2015}. Because Sentinel-2 MSI and Landsat 8/9 OLI provide complementary revisit times, their joint use as a virtual constellation coupled with machine learning is now the standard practice for mapping water quality \citep{Pahlevan2020, Cao2022}.

Although these virtual constellations work well in turbid, productive ``Case-2'' waters, they fail to transfer reliably to clear, oligotrophic high-altitude lakes due to two unresolved methodological assumptions. First, the cross-sensor harmonization coefficients required to merge Landsat and Sentinel-2 (e.g., \citealp{Roy2016}) were fitted on terrestrial targets. Because terrestrial surfaces are orders of magnitude brighter than clear water, their additive intercepts overwhelm the over-water signal. While near-infrared bias over water has been noted \citep{Pahlevan2019}, the magnitude of this mismatch and its effect on machine-learning retrievals remain unquantified. Second, to prevent spatial autocorrelation from inflating apparent accuracy, previous methods routinely block cross-validation by monitoring station \citep{Roberts2017}. However, they assume rather than test that station identity is the primary source of data leakage. In campaign-based monitoring programmes, where an entire lake is sampled within a few days under identical conditions, station-blocking may retain severe leakage and falsely inflate performance metrics. 

In this paper, we propose a rigorous validation and retrieval pipeline to overcome these challenges, using Lake Titicaca as our testing ground. Our innovation is replacing assumed methodological defaults with data-driven diagnostics. Specifically, we assemble \numMatchupsTotal{} real satellite--field match-ups and apply a classifier two-sample test (adversarial validation) \citep{LopezPaz2017} to diagnose harmonization failure. Furthermore, we introduce a null-model diagnostic to identify the true blocking unit for cross-validation. Finally, we train a Random Forest regressor under the corrected design and attach split-conformal prediction intervals \citep{Angelopoulos2023} to quantify per-prediction uncertainty. 

In contrast to previous water-quality mapping efforts \citep{AyalaIzurieta2023,Maligaya2024}, our study makes three explicit contributions:
\begin{enumerate}
\item \textbf{Quantification of Harmonization Limits:} We show that land-derived harmonization intercepts exceed the native Landsat signal by up to \numRoyRatioNir{}-fold over clear water. Consequently, a classifier can identify the source sensor with \numSensorAccAll\% accuracy from harmonized bands, proving that the model learns a sensor label rather than unified optics.
\item \textbf{Null-Model Validation Design:} We demonstrate that under campaign-blocked folds, a station-mean null model still achieves $R^2=\numNullStationRTwo$, proving that station-blocking fails to prevent leakage. We establish the campaign as the required blocking unit for this class of data. 
\item \textbf{Defensible Retrieval with Conformal Uncertainty:} Under this strict design, our model achieves an honest $R^2$ of \numHeadlineRTwo{} (RMSE \numHeadlineRmse~m). We further reveal that while split-conformal intervals are marginally calibrated (\numConformalCovNineZero\% coverage), they structurally under-cover extreme values in the clearest water, highlighting a critical limitation for operational monitoring.
\end{enumerate}

\section{Materials and methods}\label{sec:methods}

\subsection{Study area}

Lake Titicaca comprises three limnologically distinct zones
(Fig.~\ref{fig:study}): the large, deep, oligotrophic Lago Mayor; the
shallower Lago Menor (Wi\~naymarca) to the south-east; and Bah\'ia de Puno,
which receives municipal and agricultural inputs from the city of Puno and is
the most eutrophic sector. Measured transparency spans \numSecchiMin~m in
near-shore turbid water to \numSecchiMax~m in the open oligotrophic basin.

\begin{figure}[htbp]
\centering
\includegraphics[width=0.82\textwidth]{fig02_study_area.png}
\caption{Study area and monitoring network. (a) IMARPE/ALT stations with a
satellite--field match-up, coloured by mean measured Secchi depth over the
Natural Earth lake boundary; the trophic gradient is evident in the field
data alone. (b) Sampling effort per station: most stations were visited in
few campaigns, which limits the power of any station-level
analysis.}\label{fig:study}
\end{figure}

\subsection{In-situ data}

We used 881 field measurements of Secchi transparency, with co-measured
chlorophyll-$a$, suspended solids, temperature, pH, dissolved oxygen and
nutrients, from the limnological monitoring campaigns of IMARPE (Instituto
del Mar del Per\'u) on Lake Titicaca, spanning 2011--2024 across 156 stations
distributed over the whole lake. Of these, 734 include a valid Secchi
reading. The July 2019 campaign (71 stations) is published by
\citet{Siguayro2022data}; the remaining years derive from the same IMARPE
programme. As a provenance check, our 2019 zone-level means reproduce the
published values almost exactly (Lago Mayor conductivity 1513.4 vs.\
1513.5~$\mu$S~cm$^{-1}$; Lago Menor chlorophyll-$a$ 1.52 vs.\
1.52~mg~m$^{-3}$).

Three properties of this record constrain everything that follows, and we
state them here rather than in the limitations (Fig.~\ref{fig:inventory}).

First, the monitoring programme is itself almost entirely a dry-season
programme. The 2011--2024 record contains 13 campaigns with Secchi readings,
of which exactly one --- December 2012, 72 measurements --- falls in the
austral wet season (December--March). That campaign predates both sensors
used here: Landsat~8 became operational in May 2013 and Sentinel-2A in June
2015. Consequently every match-up in our dataset falls between July and
October, and this is a property of the field programme rather than of our
processing: no choice of match-up window or cloud threshold can recover
wet-season pairs that have no contemporaneous imagery. We return to what this
means for operational use in Section~\ref{sec:envelope}.

Second, sampling is episodic: campaigns exist for
\numYearFirst--\numYearLast{} with no campaigns in \numYearsMissing.

Third, the reference measurement is coarse: \numSecchiUniqueValues{} distinct
Secchi values occur in \numMatchupsSecchi{} match-ups, because disk readings
are recorded to the nearest half metre.

No synthetic or spectrally derived labels were used at any stage.



\subsection{Satellite data and cross-sensor harmonization}

Sentinel-2 MSI (harmonized surface reflectance) and Landsat 8/9 OLI
(Collection~2, Level-2; \citealp{Vermote2016}) were accessed through Google
Earth Engine \citep{Gorelick2017}. Clouds were masked with the Sentinel-2
Scene Classification Layer plus cloud probability, and with the Landsat
\texttt{QA\_PIXEL} bitmask (dilated cloud, cloud and shadow). The cirrus flag
was not applied, as it is near-permanently set at 3810~m and would mask the
entire archive. We retained the six bands common to both sensors (blue,
green, red, NIR, SWIR1, SWIR2).

Landsat surface reflectance was first harmonized to Sentinel-2 equivalent
reflectance with the per-band linear coefficients of \citet{Roy2016}, of the
form $\rho_{\mathrm{MSI}} = a + b\,\rho_{\mathrm{OLI}}$.
Section~\ref{sec:harm} evaluates whether this transform is valid over the
water of Lake Titicaca and derives an alternative. For that evaluation we
identified every in-situ event observed by both sensors within the match-up
window --- \numPairedEvents{} paired station--date observations --- and
regressed Sentinel-2 on harmonized Landsat reflectance band by band.

\subsection{Satellite--field match-ups}

For every in-situ measurement we extracted the cloud-masked median
reflectance within a 30~m buffer of the station coordinate, from imagery
acquired within $\pm10$~days. The 30~m buffer matches the coarser Landsat
pixel and averages residual geolocation error and sensor noise over several
pixels while remaining in open water, away from shoreline mixing. The
$\pm10$-day window balances the scarce cloud-free image availability at
3810~m against the slow temporal change of transparency in this large, deep
lake, and is consistent with windows used in operational $Z_{SD}$ retrieval
\citep{Kloiber2002,Zhang2022}.

This yielded \numMatchupsTotal{} match-ups (\numMatchupsSTwo{} Sentinel-2,
\numMatchupsLs{} Landsat), of which \numMatchupsSecchi{} have a Secchi value,
across \numStations{} stations and \numCampaigns{} campaign dates. Because
both sensors can observe the same field event, \numCrossSensorPairs{}
match-ups are cross-sensor pairs of a single in-situ measurement; collapsing
them leaves \numIndependentEvents{} independent events, and
Section~\ref{sec:blocking} verifies that accuracy is unchanged, confirming
the result is not an artefact of pseudo-replication.

\subsection{Features and model}

Twelve sensor-agnostic features were computed from the six common bands: the
bands themselves plus NDWI, NDTI and the ratios blue/green, green/blue,
red/green and blue/red. We trained a Random Forest regressor
\citep{Breiman2001}. Reflectance was scaled with a robust (median/IQR) scaler
fitted only on the training partition of each fold. All models were
implemented with scikit-learn \citep{Pedregosa2011}.

Model choice was settled empirically rather than assumed. We compared Random
Forest, Extremely Randomized Trees, XGBoost \citep{Chen2016} and LightGBM,
each with hyperparameters optimized by Optuna \citep{Akiba2019} inside
\emph{nested} cross-validation --- tuning on an inner campaign-blocked split
of each outer training partition, so no outer fold ever informs its own
hyperparameters --- and benchmarked the resulting gain against the spread
across random seeds. Section~\ref{sec:model} reports the outcome.

\subsection{Validation design and null models}\label{sec:valdesign}

Rather than a single split, we evaluate a hierarchy of increasingly demanding
designs \citep{Roberts2017}: random $K$-fold; GroupKFold blocked by station;
GroupKFold blocked by \emph{campaign date}; GroupKFold blocked by year; a
temporal hold-out; and leave-one-zone-out, in which an entire trophic zone is
withheld.

Crucially, we do not assume which blocking unit matters. We benchmark each
design against null models that predict, respectively, the training-set mean
of the zone, of the station and of the campaign date. A blocking unit is
doing real work if, and only if, the corresponding null model collapses under
it. This diagnostic, rather than convention, determines our primary design.

Performance is reported as out-of-fold $R^2$, RMSE and MAE, and benchmarked
against a classical log-linear blue/green band-ratio algorithm in the
standard form $\ln Z_{SD} \sim \ln(\rho_{\mathrm{blue}}/\rho_{\mathrm{green}})$
\citep{Kloiber2002}, with predictions bounded to the physically observable
range, and against a multiband linear model.

\subsection{Uncertainty and interpretation}

We attach a per-prediction interval using split-conformal prediction
\citep{Angelopoulos2023}: within each fold, a calibration subset of the
training partition yields the $(1-\alpha)$ quantile of absolute residuals,
defining a symmetric interval with finite-sample marginal coverage guarantees
\citep{Romano2019}. Because a single nominal level can be met by chance, we
report empirical coverage at 50, 80, 90 and 95\%, and also coverage
conditional on measured transparency, since marginal calibration does not
imply conditional calibration.

The model is interpreted with SHAP \citep{Lundberg2017,Lundberg2020},
computed \emph{out of fold} --- each observation explained by the model that
did not see it --- rather than on a model fitted to the whole dataset, which
inflates the apparent importance of memorized features. Trends in annual
median in-situ transparency were tested with Mann--Kendall and Sen's slope
\citep{Mann1945,Sen1968} on the full 2011--2024 in-situ record, not on the
match-up subset.

\section{Results}\label{sec:results}

\subsection{Land-derived harmonization is invalid over oligotrophic
water}\label{sec:harm}

The \citet{Roy2016} transforms carry additive intercepts of
\numRoyInterceptMin--\numRoyInterceptMax{} in reflectance units. Over Lake
Titicaca the native Landsat over-water signal is far smaller than this
(Table~\ref{tab:roy}, Fig.~\ref{fig:harm}a). In the near-infrared the median
native reflectance is effectively zero, so the intercept exceeds the signal
by a factor of \numRoyRatioNir; in SWIR1 the factor is \numRoyRatioSwirOne,
in red \numRoyRatioRed, and even in blue --- the brightest water band ---
\numRoyRatioBlue. The harmonized Landsat bands are consequently dominated by
a constant inherited from a land calibration: their medians exceed the
corresponding Sentinel-2 medians by \numLsOverSTwoGreen-fold in green and
\numLsOverSTwoNir-fold in the near-infrared.



A harmonization that succeeded would make the two sensors statistically
indistinguishable over the same water. They are not. A Random Forest
classifier trained to identify the source sensor from the harmonized features
alone reaches \numSensorAccAll\% accuracy against a \numSensorAccChance\%
majority-class baseline (Fig.~\ref{fig:harm}b); restricting it to the visible
bands gives \numSensorAccVisible\%, and to scale-invariant band ratios
\numSensorAccRatios\%. This is a classifier two-sample test
\citep{Friedman2004,LopezPaz2017}: if two samples come from the same
distribution, no classifier can separate them above chance. Sensor identity
therefore remains fully encoded in the harmonized reflectance, and any model
trained on the pooled archive can read it. The test is not new, but to our
knowledge it has not been used as a harmonization acceptance check; it costs
one model fit and it is decisive.

Refitting locally on the \numPairedEvents{} station--date events observed by
both sensors gives water-specific coefficients (Table~\ref{tab:localcoef}).
These restore the spectral signature almost exactly --- the median
Landsat/Sentinel-2 ratio falls from 15.4 to about 1.0 across all six bands
(Fig.~\ref{fig:harm}d) --- but per-observation agreement remains uneven
(Fig.~\ref{fig:harm}c). Green recovers usable agreement
($r=\numLocalRGreen$), and NIR and SWIR reach $r\approx0.6$. Blue does not:
$r=\numLocalRBlue$. This is the most consequential result of the section,
because SHAP attributes \numShapShareVerdeAzul\% of the retrieval to
green/blue bands and ratios (Section~\ref{sec:shap}). The band on which the
retrieval most depends is the band on which the two sensors least agree.



\begin{figure}[htbp]
\centering
\includegraphics[width=\textwidth]{fig01_harmonization_failure.png}
\caption{Why land-derived harmonization fails over clear water. (a) The Roy
intercept relative to the native over-water Landsat signal, by band. (b)
Accuracy of a classifier predicting the source sensor from harmonized
features; under a successful harmonization it would sit at chance. (c)
Sentinel-2 versus harmonized Landsat reflectance on the \numPairedEvents{}
paired station--date events; blue agrees worst. (d) Median spectral signature
before and after local water recalibration; labels are the
Landsat/Sentinel-2 ratio.}\label{fig:harm}
\end{figure}

\subsection{The campaign, not the station, is the blocking
unit}\label{sec:blocking}

Table~\ref{tab:blocking} reports the retrieval under each validation design
alongside the null models (Fig.~\ref{fig:val}). Blocking by station changes
nothing relative to a random split (both $R^2=\numRandomRTwo$), which under
the usual reading would be taken as evidence of no spatial leakage. The null
models show why that reading is wrong: under campaign-blocked folds the
station-mean null still attains $R^2=\numNullStationRTwo$ and the zone-mean
null $R^2=\numNullZonaRTwo$, whereas the campaign-date-mean null collapses to
$\numNullCampaignDateRTwo$. Dependence in this dataset lives between
campaigns, not between stations, because a single campaign samples the whole
lake within a few days under one set of atmospheric and limnological
conditions. Station-blocked folds leave that structure intact.

Under campaign-blocked validation --- our primary design --- the Random
Forest attains $R^2=\numHeadlineRTwo$ and RMSE \numHeadlineRmse~m over
\numHeadlineN{} match-ups, essentially unbiased (\numHeadlineBias~m).
Blocking by year is stricter still ($R^2=\numYearRTwo$). Collapsing
cross-sensor pseudo-replicates to \numOneEventN{} independent events leaves
accuracy unchanged ($R^2=\numOneEventRTwo$).

\begin{table}[htbp]
\caption{Validation design, null models and blocking units (Random Forest,
$n=812$ unless noted). A blocking unit is effective only if the corresponding
null model collapses under it.}\label{tab:blocking}
\tabfit
\centering
\resizebox{\linewidth}{!}{\begin{tabular}{@{}llrr@{}}
\toprule
& Setting & $R^2$ & RMSE (m) \\
\midrule
\multirow{3}{*}{Null models} & Zone-mean            & 0.287  & 2.40 \\
                             & Station-mean         & 0.208  & 2.53 \\
                             & Campaign-date-mean   & $-0.012$ & 2.86 \\
\midrule
\multirow{5}{*}{Blocking unit} & Random $K$-fold        & 0.596 & 1.81 \\
                               & By station             & 0.596 & 1.81 \\
                               & \textbf{By campaign date} & \textbf{0.574} & \textbf{1.86} \\
                               & By year                & 0.495 & 2.02 \\
                               & One record per event ($n=540$) & 0.571 & 1.77 \\
\midrule
\multirow{3}{*}{Zone extrapolation} & Leave out Lago Menor ($n=133$)  & 0.247   & 1.90 \\
                                    & Leave out Bah\'ia de Puno ($n=88$) & 0.036 & 2.11 \\
                                    & Leave out Lago Mayor ($n=591$)  & $-0.047$ & 2.52 \\
\midrule
Temporal & Train $\le$2019, test 2022 \& 2024 ($n=226$) & 0.507 & 2.63 \\
\bottomrule
\end{tabular}}
\end{table}



\subsection{Retrieval accuracy and benchmarks}

Against this design the classical blue/green band ratio attains
$R^2=\numClassicalBoundedRTwo$ (RMSE \numClassicalBoundedRmse~m) and a
multiband linear model $R^2=\numLinearRTwo$ (RMSE \numLinearRmse~m)
(Table~\ref{tab:bench}, Fig.~\ref{fig:retr}). We report one methodological
caution alongside these figures. An unbounded implementation of the same
band-ratio algorithm --- log-transformed and back-transformed without
constraining predictions to the observable range --- produces a single
prediction of \numClassicalBrokenMaxpred~m against an observed maximum of
\numSecchiMax~m, which alone drives $R^2$ to \numClassicalBrokenRTwo{} and
RMSE to \numClassicalBrokenRmse~m (Fig.~\ref{fig:retr}d). Benchmarks of
classical algorithms reported without physical bounds should be treated with
caution: a single unconstrained extrapolation manufactures an arbitrarily
large advantage for the machine-learning model.



\begin{figure}[htbp]
\centering
\includegraphics[width=\textwidth]{fig04_retrieval_and_benchmark.png}
\caption{Retrieval under the primary design. (a) Out-of-fold predicted versus
measured Secchi depth. (b) Residuals by trophic zone; the retrieval
overestimates in eutrophic Bah\'ia de Puno. (c) Benchmark against classical
and linear baselines. (d) The unbounded classical implementation in log
scale: one impossible prediction destroys the metric.}\label{fig:retr}
\end{figure}

\subsection{Model family and tuning do not change the
picture}\label{sec:model}

We compared four learner families, each tuned by Optuna inside nested
cross-validation, against target transformations and reduced feature sets
(Table~\ref{tab:model}). The best configuration --- an equal-weight ensemble
of the four tuned models --- reached $R^2=\numBestModelRTwo$, an improvement
of \numTuningGainRTwo{} over the untuned Random Forest. That gain does exceed
the spread across random seeds (\numSeedNoiseRange{} in $R^2$), so it is not
noise; but it corresponds to 0.02~m of RMSE on a lake whose reference
measurement is recorded to the nearest half metre, and it changes no
conclusion in this article. Neither log-transforming the target nor tuning
individual learners helped more than that, and restricting the feature set
cost accuracy.

We therefore retain the untuned Random Forest as the reported model: it is
simpler, faster and easier to reproduce, and reporting a tuned ensemble for a
0.009 gain would overstate the precision of the product. The ceiling here is
set by the data --- a dry-season-only record, a reference measurement
quantized at half a metre, and a target whose variance is dominated by a
three-level trophic gradient --- not by the choice of learner.



\subsection{Where the retrieval stops working}

Withholding an entire trophic zone collapses the retrieval
(Table~\ref{tab:blocking}): $R^2$ falls to \numZoneoutMenorRTwo{} for Lago
Menor, \numZoneoutPunoRTwo{} for Bah\'ia de Puno and \numZoneoutMayorRTwo{}
for Lago Mayor. The model interpolates within optical regimes it has seen and
does not extrapolate to a water type absent from training. On the temporal
hold-out (train $\le$\numTemporalCut, test 2022 and 2024) it retains
$R^2=\numTemporalRTwo$, though \numTemporalShared{} of the
\numTemporalTestStations{} test stations also appear in training
(\numTemporalOverlapPct\%), so this measures temporal rather than fully
independent generalization.

Error is unevenly distributed (Table~\ref{tab:zone}). The retrieval carries a
positive bias of \numZoneBahiaDePunoBias~m in Bah\'ia de Puno, the most
eutrophic and most vulnerable zone, because the clear-water majority pulls
predictions upward.



Two features of the harmonized archive should be read together with these
numbers. Removing the two bands most corrupted by the land-derived intercepts
lowers $R^2$ only from \numHeadlineRTwo{} to \numVisibleRTwo; using
scale-invariant ratios alone gives \numRatiosRTwo; and applying our
water-specific recalibration gives \numLocalReharmRTwo. Training on each
sensor separately gives \numSTwoOnlyRTwo{} (Sentinel-2, $n=\numSTwoOnlyN$)
and \numLsOnlyRTwo{} (Landsat, $n=\numLsOnlyN$). The pooled model's small
advantage over these variants is consistent with it exploiting sensor
identity as an additional covariate rather than with genuine cross-sensor
synergy.

\subsection{Calibrated per-prediction uncertainty}

Split-conformal intervals are well calibrated \emph{marginally} at every
nominal level tested: empirical coverage is \numConformalCovFiveZero\% at
50\% nominal, \numConformalCovEightZero\% at 80\%,
\numConformalCovNineZero\% at 90\% and \numConformalCovNineFive\% at 95\%,
a maximum calibration error of \numConformalMaxError{} percentage points
(Fig.~\ref{fig:unc}a).

They are not calibrated \emph{conditionally}, and this matters more for
operational use than the marginal result. Coverage falls to
\numConformalCovClearest\% in the clearest water ($Z_{SD}>12.5$~m) and
\numConformalCovTurbidest\% in the most turbid ($Z_{SD}<5$~m), while
exceeding the nominal level in the middle of the range
(Fig.~\ref{fig:unc}c). The mechanism is structural rather than incidental:
split conformal derives a single global quantile of absolute residuals, so
the interval has constant width (\numConformalWidthNineZero~m at 90\%)
everywhere and cannot adapt to the local difficulty of the retrieval. It
therefore buys its marginal guarantee by over-covering in the well-sampled
middle of the range and under-covering at both extremes --- which are
precisely where management decisions are made. A locally adaptive variant
such as conformalized quantile regression \citep{Romano2019} would be
required for reliable conditional coverage, and we recommend reporting
coverage conditional on the predicted value as standard practice whenever
conformal intervals accompany an operational product.

For a lake whose transparency spans \numSecchiMin--\numSecchiMax~m, a
$\pm3$~m interval is informative for zone-scale contrasts and multi-metre
change, and is not informative for sub-metre differences.



\subsection{Interpretation and retrievability}\label{sec:shap}

Out-of-fold SHAP attributes \numShapShareVerdeAzul\% of the retrieval to
green and blue bands and their ratios, \numShapShareNirSwir\% to NIR and
SWIR, and the remainder to other indices (Fig.~\ref{fig:shap}a). Green/blue
alone accounts for 37\%. The dependence is monotone and of the sign expected
from water-clarity optics (Fig.~\ref{fig:shap}b): a higher green-to-blue
ratio pushes the prediction toward less transparent water. The retrieval is
therefore grounded in optics rather than location memorization --- while
resting on the band pair with the worst cross-sensor agreement.

With real labels and the same campaign-blocked design, no other variable is
retrievable (Fig.~\ref{fig:shap}c): chlorophyll-$a$ reaches
$R^2=\numRetrChlRTwo$, suspended solids \numRetrTssRTwo, in-situ water
temperature \numRetrTempInsituRTwo, and Landsat thermal surface temperature
\numRetrLstThermalRTwo{} ($r=0.14$, $n=\numRetrLstThermalN$). Their in-situ
ranges --- chlorophyll-$a$ 0--4.6~mg~m$^{-3}$, suspended solids mostly below
10~mg~L$^{-1}$ --- lie below the optical and thermal detection limits of
these sensors. We subject the negative claim to the same design as the
positive one, and report it to caution against applying chlorophyll
algorithms \citep{Gitelson2009,Mishra2012} or thermal retrievals
\citep{RuizVerdu2016} in clear high-altitude water without local validation.

\begin{figure}[htbp]
\centering
\includegraphics[width=\textwidth]{fig08_shap_and_retrievability.png}
\caption{Interpretation and operational envelope. (a) Out-of-fold SHAP
importance; green and blue dominate. (b) SHAP dependence on the leading
feature, by trophic zone. (c) Retrievability of each water-quality variable
under the same campaign-blocked design; only transparency clears a useful
threshold.}\label{fig:shap}
\end{figure}

\subsection{Temporal patterns and the fragility of trend
claims}\label{sec:trends}

Mann--Kendall tests on annual median in-situ transparency over the full
2011--2024 record show no statistically significant trend in any zone
(Bah\'ia de Puno: Sen slope \numTrendBahiaDePunoSlope~m~yr$^{-1}$,
$p=\numTrendBahiaDePunoP$; Lago Mayor: \numTrendLagoMayorSlope~m~yr$^{-1}$,
$p=\numTrendLagoMayorP$; Fig.~\ref{fig:trends}a).

That conclusion is fragile in a way worth reporting explicitly. Restricting
the series to 2013 onwards --- the natural choice if trends were computed on
the satellite match-up subset, whose record begins with Landsat~8 --- turns
both Bah\'ia de Puno ($p=0.017$) and Lago Mayor (Sen $+0.42$~m~yr$^{-1}$,
$p=0.012$) significantly positive. Across start years from 2011 to 2014,
\numTrendFlips{} of \numTrendCombinations{} zone--start-year combinations
change verdict (Fig.~\ref{fig:trends}b). The year 2013 has the lowest annual
medians on record, so beginning there manufactures an upward trend. With
9--11 annual values and three gap years the test has low power, and the
correct reading is absence of evidence, not evidence of stability. We
recommend that trend analyses in campaign-based lake records report this
sensitivity as a matter of course.



\section{Discussion}\label{sec:discussion}

\subsection{Implications for harmonized multi-mission archives}

The virtual-constellation approach \citep{Claverie2018,Pahlevan2019} has
become the default for multi-decadal lake monitoring, and the
\citet{Roy2016} coefficients are its most reused transform. Our results place
a boundary on where that default is safe. Those coefficients were fitted on
terrestrial targets and carry additive intercepts appropriate to surfaces one
to two orders of magnitude brighter than clear water. Over Lake Titicaca the
intercept exceeds the native over-water signal by a factor of
\numRoyRatioNir{} in the near-infrared and by \numRoyRatioBlue{} even in
blue, and the two sensors remain \numSensorAccAll\% separable after the
transform. That water reflectance from these products is unreliable is
already documented \citep{Pahlevan2019,HLSguide}; what our paired
observations add is the magnitude --- the problem is not confined to the
near-infrared, and it is large enough to make the harmonized bands function
as sensor labels for a machine-learning model --- and the finding that local
recalibration repairs the aggregate spectral signature without repairing
per-observation agreement in blue.

The practical consequence is direct. A machine-learning model trained on a
pooled, land-harmonized archive can achieve apparently good accuracy while
having learned two sensor-specific mappings joined by a readable sensor
label. That is not a harmonized retrieval, and it will not behave as one when
applied to a scene from either sensor alone. The diagnostic is cheap and we
recommend it as routine: train a classifier to predict the source sensor from
the harmonized features and report its accuracy against the majority-class
baseline. If the sensors remain separable, the harmonization has not done its
job. For clear oligotrophic lakes we further recommend deriving
water-specific coefficients from same-day paired observations, reporting
per-band agreement, and --- where blue-band agreement is as poor as we find
here --- reporting single-sensor retrievals alongside any pooled product.

\subsection{Implications for validation design}

Blocking by station is the standard remedy for spatial autocorrelation in
this literature, and in our data it changes nothing: random and
station-blocked folds give identical accuracy. Read conventionally, that
would be reassuring. The null models show it is not. Under campaign-blocked
folds the station-mean null still explains 21\% of variance and the zone-mean
null 29\%, while the campaign-date-mean null explains none. In a
campaign-based monitoring programme the whole lake is sampled within a few
days under shared atmospheric and limnological conditions, so observations
cluster by campaign far more strongly than by location; station-blocked folds
leave that structure entirely intact.

The general point is that the blocking unit should be diagnosed, not assumed.
Reporting null models under each candidate design costs little and identifies
the unit that carries the dependence. This applies well beyond Lake Titicaca:
most large-lake monitoring programmes sample in campaigns, so the same
mismatch between the conventional blocking unit and the real dependence
structure is likely to be widespread.

For our retrieval, the honest figure is $R^2=\numHeadlineRTwo{}$ (RMSE
\numHeadlineRmse~m) under campaign-blocked validation, against
\numNullZonaRTwo{} for a zone-mean null: the model roughly doubles the skill
available from knowing only which basin a sample came from. That is a real
but modest contribution, and stating it that way is more useful to a manager
than an unqualified $R^2$.

\subsection{Operational envelope for Lake Titicaca}\label{sec:envelope}

Our RMSE of \numHeadlineRmse~m spans roughly 12\% of the observed $Z_{SD}$
range and sits below the published range for $Z_{SD}$ retrieval
($R^2\approx0.6$--0.9; \citealp{Zhang2022,Lee2015}), obtained here under a
predominantly clear oligotrophic regime and a stricter validation design. It
is also approaching the granularity of the reference measurement itself:
disk readings are recorded to the nearest half metre, giving only
\numSecchiUniqueValues{} distinct values across \numMatchupsSecchi{}
match-ups.

Four constraints bound operational use. The retrieval carries a
\numZoneBahiaDePunoBias~m positive bias in Bah\'ia de Puno --- the most
eutrophic and most vulnerable zone --- so mapped values there must be read
with their intervals and confirmed by field sampling. Leave-one-zone-out
results show the model cannot extrapolate to an unsampled optical regime;
integration must be incremental, each new campaign both validating the
current map and extending the training set. The retrieval is calibrated and validated for
the austral dry season only, and the reason is worth stating precisely
because it determines the remedy. It is not that wet-season imagery is
unusable at 3810~m, nor that our match-up criteria excluded it: it is that
the field programme has produced a single wet-season campaign in fourteen
years, and that campaign predates both satellites. Wet-season blooms in
Bah\'ia de Puno --- the phenomenon a transparency product would most usefully
warn about --- are therefore entirely unrepresented, and the model should not
be applied between December and March.

This converts the limitation into a specific, actionable recommendation for
the monitoring authorities. Extending the protocol to the wet season does not
require a better algorithm; it requires field data that do not yet exist. A
single additional wet-season campaign per year, timed to any cloud-free
window and covering the three trophic zones, would within a few years supply
enough independent campaign-blocks to calibrate and validate a wet-season
retrieval under the same design used here. Because the binding constraint is
sampling rather than sensing, this is among the cheapest available
improvements to Lake Titicaca's water-quality monitoring capability. And because conformal coverage degrades at both
extremes of the transparency range, the nominal 90\% interval should be
treated as optimistic in the clearest and most turbid water.

Within those bounds the product is useful to the regional monitoring
authorities (IMARPE, ALT, OEFA, and the binational management framework). A
calibrated $\pm3$~m interval resolves zone-scale contrasts and multi-metre
change, which is the scale at which clarity loss in Bah\'ia de Puno would be
actionable, and a single transboundary product under one calibration supports
comparable assessment across the Peru--Bolivia boundary where coordinated
field monitoring is logistically difficult. It does not resolve sub-metre
change and should not be presented as doing so.

\subsection{Limitations}

Beyond the constraints already stated, three limitations bound this work. The
in-situ record is episodic, with no campaigns in \numYearsMissing{} and only
one wet-season campaign in fourteen years, so the temporal hold-out rests on
two test years, the trend analysis on 9--11 annual values, and the seasonal
envelope on none. Transparency above approximately
15~m saturates the optical signal. And this is a single lake: the two
diagnostics we propose are general, but the water-specific coefficients in
Table~\ref{tab:localcoef} are local to Titicaca and must be refitted
elsewhere.

\section{Conclusions}\label{sec:conclusions}

Using \numMatchupsTotal{} satellite--field match-ups over twelve years, we
show that two practices imported from the turbid-water remote-sensing
literature fail in a clear high-altitude large lake. Cross-sensor
harmonization coefficients derived over land leave Sentinel-2 and Landsat
\numSensorAccAll\% separable over oligotrophic water, because their additive
intercepts exceed the native over-water signal by up to a factor of
\numRoyRatioNir; and validation blocked by monitoring station leaves the
dominant dependence --- between campaigns --- entirely intact.

Correcting both, a Random Forest retrieves Secchi depth with
$R^2=\numHeadlineRTwo$ and RMSE \numHeadlineRmse~m, roughly double the skill
of a zone-mean null, with conformal prediction intervals that are calibrated
marginally but under-cover at the extremes of the transparency range. Skill
collapses under extrapolation to an unsampled trophic zone;
chlorophyll-$a$, suspended solids and thermal temperature are not retrievable
at all; and no transparency trend is significant once the full in-situ record
is used rather than the satellite-era subset.

We therefore recommend two routine checks for multi-mission lake
water-quality products, neither of which requires new methodology: report the
accuracy of a classifier two-sample test on the harmonized features, and
report null models under each candidate validation blocking unit. Both cost
one model fit, and both distinguish genuine optical retrieval from sensor
labelling and sampling design. For Lake
Titicaca the resulting product is a dry-season, zone-scale transparency
indicator with honest $\pm3$~m intervals --- less than has sometimes been
claimed for such lakes, and more than the region has had.

\appendix
\section{Supplementary Figures and Tables}

\begin{figure}[htbp]
\centering
\includegraphics[width=\textwidth]{fig05_data_inventory.png}
\caption{What the dataset actually contains. (a) Match-ups by year and month;
three years have no campaign at all, so a hold-out described as
``2022--2024'' contains only 2022 and 2024. (b) Seasonal coverage: zero
match-ups in the wet season, when blooms occur in Bah\'ia de Puno. (c) Secchi
distribution by trophic zone; the reference measurement is quantized at
roughly half a metre.}\label{fig:inventory}
\end{figure}

\begin{table}[htbp]
\caption{Effect of the \citet{Roy2016} OLI$\rightarrow$MSI transform over
Lake Titicaca water. Native OLI is the Collection-2 surface reflectance
recovered by inverting the transform. Reflectance is dimensionless.}
\label{tab:roy}
\tabfit
\centering
\resizebox{\linewidth}{!}{\begin{tabular}{@{}lrrrrr@{}}
\toprule
Band & Roy intercept $a$ & Native OLI & Intercept / & Harmonized & Harmonized / \\
     &                   & median     & signal      & LS median  & S2 median \\
\midrule
Blue   & 0.0183 & $+0.0072$ & 2.5  & 0.0247 & 1.74 \\
Green  & 0.0123 & $+0.0089$ & 1.4  & 0.0206 & 1.61 \\
Red    & 0.0123 & $+0.0013$ & 10.6 & 0.0134 & 2.66 \\
NIR    & 0.0448 & $-0.0006$ & 69.8 & 0.0443 & 15.42 \\
SWIR1  & 0.0306 & $+0.0013$ & 23.2 & 0.0317 & 13.34 \\
SWIR2  & 0.0116 & $+0.0013$ & 8.2  & 0.0129 & 5.70 \\
\bottomrule
\end{tabular}}
\end{table}

\begin{table}[htbp]
\caption{Water-specific LS$\rightarrow$S2 recalibration fitted on
\numPairedEvents{} paired station--date observations, of the form
$\rho_{\mathrm{S2}} = a + b\,\rho_{\mathrm{LS,Roy}}$.}\label{tab:localcoef}
\tabfit
\centering
\resizebox{\linewidth}{!}{\begin{tabular}{@{}lrrrl@{}}
\toprule
Band & $a$ & $b$ & Pearson $r$ & Agreement \\
\midrule
Blue   & $+0.0084$ & 0.258 & 0.26 & unusable \\
Green  & $-0.0009$ & 0.665 & 0.71 & good \\
Red    & $-0.0005$ & 0.427 & 0.52 & weak \\
NIR    & $-0.0250$ & 0.631 & 0.62 & acceptable \\
SWIR1  & $-0.0223$ & 0.770 & 0.64 & acceptable \\
SWIR2  & $-0.0057$ & 0.615 & 0.60 & acceptable \\
\bottomrule
\end{tabular}}
\end{table}

\begin{figure}[htbp]
\centering
\includegraphics[width=0.85\textwidth]{fig03_validation_design.png}
\caption{Diagnosing the blocking unit. (a) Null models evaluated under
campaign-blocked folds: only the campaign-mean null collapses, so only
campaign blocking is doing work. (b) Out-of-fold $R^2$ under each validation
design and under extrapolation to an unseen trophic zone. That station
blocking matches the random split is not evidence of no leakage; it is
evidence that station blocking retains nothing.}\label{fig:val}
\end{figure}

\begin{table}[htbp]
\caption{Head-to-head benchmark, campaign-blocked GroupKFold, $n=812$.}
\label{tab:bench}
\tabfit
\centering
\resizebox{\linewidth}{!}{\begin{tabular}{@{}lrrrr@{}}
\toprule
Algorithm & $R^2$ & RMSE & MAE & Bias \\
\midrule
Classical blue/green, unbounded$^{a}$      & $-2.458$ & 5.29 & 1.91 & $-0.19$ \\
Classical blue/green, bounded$^{b}$        & 0.097 & 2.71 & 2.13 & $-0.46$ \\
Multiband linear                           & 0.323 & 2.34 & 1.78 & $-0.01$ \\
Random Forest (this study)                 & \textbf{0.574} & \textbf{1.86} & \textbf{1.44} & $-0.02$ \\
\bottomrule
\multicolumn{5}{@{}l}{\footnotesize RMSE, MAE and bias in metres.}\\
\multicolumn{5}{@{}l}{\footnotesize $^{a}$ back-transformed without bounding predictions; see text.}\\
\multicolumn{5}{@{}l}{\footnotesize $^{b}$ standard form of \citet{Kloiber2002}.}\\
\end{tabular}}
\end{table}

\begin{table}[htbp]
\caption{Model selection under nested cross-validation (Optuna inside an
inner campaign-blocked split of each outer training partition). Selected
configurations; the full grid is in the repository.}\label{tab:model}
\tabfit
\centering
\resizebox{\linewidth}{!}{\begin{tabular}{@{}llrr@{}}
\toprule
Model & Features & $R^2$ & RMSE (m) \\
\midrule
Ensemble of the four tuned models  & all 12 & 0.583 & 1.84 \\
LightGBM, tuned                    & all 12 & 0.581 & 1.84 \\
XGBoost, tuned                     & all 12 & 0.577 & 1.85 \\
Random Forest, tuned               & all 12 & 0.576 & 1.85 \\
\textbf{Random Forest, untuned (reported)} & all 12 & \textbf{0.574} & \textbf{1.86} \\
Extremely Randomized Trees, untuned & all 12 & 0.573 & 1.86 \\
XGBoost, tuned, $\log Z_{SD}$      & all 12 & 0.561 & 1.89 \\
LightGBM, tuned                    & visible only & 0.544 & 1.92 \\
LightGBM, tuned                    & ratios only  & 0.488 & 2.04 \\
\midrule
\multicolumn{4}{@{}l}{\footnotesize Seed-to-seed spread of the untuned Random Forest: $R^2$ range 0.003 over five seeds.}\\
\bottomrule
\end{tabular}}
\end{table}

\begin{table}[htbp]
\caption{Out-of-fold error by trophic zone, campaign-blocked validation.}
\label{tab:zone}
\tabfit
\centering
\resizebox{\linewidth}{!}{\begin{tabular}{@{}lrrrr@{}}
\toprule
Zone & $R^2$ & RMSE (m) & Bias (m) & $n$ \\
\midrule
Bah\'ia de Puno & 0.221 & 1.90 & $+1.01$ & 88  \\
Lago Menor      & 0.211 & 1.94 & $+0.08$ & 133 \\
Lago Mayor      & 0.448 & 1.83 & $-0.19$ & 591 \\
\bottomrule
\end{tabular}}
\end{table}

\begin{figure}[htbp]
\centering
\includegraphics[width=\textwidth]{fig07_conformal_uncertainty.png}
\caption{Conformal prediction intervals under campaign-blocked validation.
(a) Empirical versus nominal coverage: marginally calibrated across the
50--95\% range. (b) The 90\% interval against the measured value. (c)
Coverage conditional on measured transparency. Because split conformal uses a
single global residual quantile, the interval width is constant
($\approx$\numConformalWidthNineZero~m) and cannot adapt, so the marginal
guarantee is bought by over-covering in the middle of the range and
under-covering at both extremes.}\label{fig:unc}
\end{figure}

\begin{figure}[htbp]
\centering
\includegraphics[width=\textwidth]{fig06_trends_and_start_year.png}
\caption{Transparency trends and their fragility. (a) Annual median in-situ
transparency by zone over the full 2011--2024 record, with Mann--Kendall/Sen
statistics; no trend is significant. (b) The same test with different start
years: beginning the series in 2013, the lowest-median year on record, turns
two zones significantly positive.}\label{fig:trends}
\end{figure}

\section*{CRediT authorship contribution statement}

\textbf{Dina Maribel Yana-Yucra:} Conceptualization, Methodology,
Investigation, Data curation, Writing -- original draft.
\textbf{Richar Andre Vilca-Solorzano:} Methodology, Software, Formal
analysis, Validation, Visualization, Writing -- original draft, Writing --
review \& editing.
\textbf{Fred Torres-Cruz:} Supervision, Resources, Writing -- review \&
editing.

\section*{Declaration of competing interest}

The authors declare that they have no known competing financial interests or
personal relationships that could have appeared to influence the work
reported in this paper.

\section*{Data availability}

\noindent
\begin{minipage}[c]{0.74\textwidth}
\raggedright
The satellite--field match-up datasets and the complete analysis pipeline,
including the scripts that generate every number, table and figure in this
article, are openly available at:

\medskip
{\small\sloppy\url{https://github.com/Andre031222/harmonized-titicaca-transparency}\par}

\medskip
Raw imagery is freely available via Google Earth Engine; the in-situ
measurements derive from published IMARPE/ALT and OEFA monitoring records.
\end{minipage}\hfill
\begin{minipage}[c]{0.20\textwidth}
\centering

\includegraphics[width=0.85\linewidth]{qr_repo.png}\\[3pt]
{\footnotesize Repository}
\end{minipage}

\section*{Acknowledgements}

We thank IMARPE, ALT and OEFA for the publicly available monitoring records
that made this study possible. This research received no external funding.

\bibliographystyle{elsarticle-harv}
\bibliography{references}

\end{document}
